\section{Methods}\label{sec:methods}

\subsection{Gyroid Geometry}\label{sec:geom}

The gyroid surface is defined by Eq.~\eqref{eq:gyroid}.
The solid phase occupies the region $\varphi(\mathbf{x}) \ge t$,
where $t \in (-0.5,\; 0.5)$ controls wall thickness and porosity
$\varepsilon$.  The unit-cell size $a$ determines the characteristic
pore dimension.  For each $(a, t)$ pair, the voxelized domain is
constructed on a Cartesian grid with resolution $\Delta x = \SI{0.2}{\milli\meter}$,
yielding $N_x = N_y = 131$ and $N_z = 2a/\Delta x$.

The specific surface area $S_v$ (surface area per unit volume, \si{\per\meter})
and porosity $\varepsilon$ are computed directly from the voxelized geometry.

\subsection{Lattice Boltzmann Solver}\label{sec:lbm}

A D3Q19 MRT-LBM scheme is employed~\cite{ref_mrt_lbm}.
The solver is implemented in Taichi (v1.7.4) with CUDA backend.
Key simulation parameters are listed in Table~\ref{tab:lbm_params}.

\begin{table}[htbp]
\caption{LBM simulation parameters.}\label{tab:lbm_params}
\centering
\begin{tabular}{@{}lll@{}}
\toprule
Parameter & Symbol & Value \\
\midrule
Lattice & --- & D3Q19 \\
Collision model & --- & MRT \\
Grid resolution & $\Delta x$ & \SI{0.2}{\milli\meter} \\
Domain (cross-section) & $N_x \times N_y$ & $131 \times 131$ \\
Domain (streamwise) & $N_z$ & $2a/\Delta x$ \\
Boundary condition & --- & Periodic + body force \\
Framework & --- & Taichi 1.7.4, \texttt{arch=cuda} \\
\bottomrule
\end{tabular}
\end{table}

The no-slip condition at solid walls is enforced via the bounce-back scheme.
Periodic boundary conditions are applied in the streamwise ($z$) direction,
and a uniform body force $g_\mathrm{lbm}$ drives the flow.
Darcy permeability is obtained from the steady-state velocity field:
\begin{equation}\label{eq:darcy}
  K = \frac{\mu \, \langle u_z \rangle}{(\Delta P / L)},
\end{equation}
where $\mu$ is the dynamic viscosity, $\langle u_z \rangle$ the
volume-averaged streamwise velocity, and $L$ the domain length.

\subsection{Bayesian Optimization}\label{sec:bo}

The BO framework uses Gaussian process regression with the
Lower Confidence Bound (LCB) acquisition function
($\kappa = 3$)~\cite{ref_bo_lcb}, implemented via
\texttt{scikit-optimize}.  The design space is:

\begin{table}[htbp]
\caption{Bayesian optimization configuration.}\label{tab:bo_config}
\centering
\begin{tabular}{@{}ll@{}}
\toprule
Parameter & Value \\
\midrule
Design variable $a$ & $[\num{3.0},\; \num{12.0}]$ \si{\milli\meter} \\
Design variable $t$ & $(-0.5,\; 0.5)$ \\
Acquisition function & LCB, $\kappa = 3$ \\
Initial samples & 20 (random) \\
Iterations & 80 \\
Total evaluations & 100 \\
\bottomrule
\end{tabular}
\end{table}

The objective is to minimize $\Delta P_\mathrm{Darcy}$
while simultaneously maximizing $S_v$.
A design is marked \emph{feasible} if the simulation converges
and satisfies physical consistency checks.

\subsection{Post-processing}\label{sec:postproc}

\paragraph{Forchheimer analysis.}
For the top-5 designs, five body-force levels
($g_\mathrm{lbm} \in \{10^{-6},\; 5{\times}10^{-6},\; 5{\times}10^{-5},\;
5{\times}10^{-4},\; 2{\times}10^{-3}\}$) are applied.
The extended Darcy--Forchheimer equation,
\begin{equation}\label{eq:forch}
  \frac{\Delta P / L}{u} = \frac{\mu}{K} + \beta\,\rho\,u,
\end{equation}
is fitted via linear regression of $(\Delta P/L)/u$ versus $u$
to extract permeability $K$ and inertial coefficient $\beta$.

\paragraph{Flow metrics.}
Mixing ratio ($v_\mathrm{trans}/|v_z|$), velocity uniformity
($\sigma_{v_z}$), and hydraulic tortuosity $\tau$ are evaluated
at the mid-plane cross-section.

\paragraph{GHSV sweep.}
Pressure drop at five gas hourly space velocities
(GHSV $= 5000$--$\num{60000}\;\si{\per\hour}$) is computed
from the Darcy relation $\Delta P = u_\mathrm{in}\,\mu\,L_\mathrm{cat}/K$.
